\documentclass[11pt,a4paper]{article}

% --- Paquetes ---
\usepackage[utf8]{inputenc}
\usepackage[T1]{fontenc}
\usepackage[spanish]{babel}
\usepackage{amsmath, amssymb, amsthm}
\usepackage{geometry}
\usepackage{tikz}
\usepackage{pgfplots}
\usepackage{enumitem}
\usepackage{hyperref}
\usepackage{xcolor}
\usepackage{tcolorbox}

\geometry{margin=2.5cm}
\pgfplotsset{compat=1.18}

% --- Definiciones ---
\newcommand{\barvar}[1]{\overline{#1}}

\newtcolorbox{definicionbox}[1][]{colback=yellow!10, colframe=red!70!black, title=#1, fonttitle=\bfseries}
\newtcolorbox{resultadobox}[1][]{colback=yellow!10, colframe=red!80!black, title=#1, fonttitle=\bfseries}
\newtcolorbox{explicacionbox}[1][]{colback=red!5, colframe=red!60!black, title=#1, fonttitle=\bfseries}
\newtcolorbox{notabox}[1][]{colback=pink!10, colframe=red!50!black, title=#1, fonttitle=\bfseries}

% --- Documento ---
\title{\textbf{Clase 5: Teoría de la Preferencia por la Liquidez y Equilibrio IS-LM} \\ \large Macroeconomía II}
\author{Universidad Católica del Norte}
\date{07 de septiembre, 2026}

\begin{document}
\maketitle
\tableofcontents
\newpage

% ============================================================
\section{Teoría de la Preferencia por la Liquidez}
% ============================================================

\begin{definicionbox}[Curva LM]
La Curva LM muestra la relación entre el tipo de interés y el nivel de producción donde el \textbf{mercado del dinero} está en equilibrio. Todo empieza por la teoría de la preferencia por la liquidez.
\end{definicionbox}

\subsection{Nominal vs Real}

\begin{notabox}[Demanda Nominal vs Real de Dinero]
\textbf{Demanda Nominal de Dinero:} Demanda que hace un individuo de cierta suma de dinero.

\textbf{Demanda Real de Dinero:} Se expresa en términos del poder adquisitivo del dinero, es decir, el número de unidades de bienes que compra el dinero.
\end{notabox}

\subsubsection{Ejemplo}

Demanda por dinero nominal $= \$100.000$

Para calcular la demanda real, necesitamos el nivel de precios. Usamos el \textbf{Índice de Precios al Consumidor (IPC)}, que mide el costo promedio de una canasta de bienes y servicios representativa de una familia promedio.

\begin{itemize}
    \item IPC Julio 2023 $\rightarrow$ 132,20 puntos del IPC (Base = 100, año 2018).
    \item Nivel de Precios en términos relativos: $\dfrac{132{,}20}{100} = 1{,}322$
\end{itemize}

\[
\text{Demanda Real de Dinero} = \frac{\text{Demanda Nominal}}{\text{Nivel de Precios}} = \frac{\$100.000}{1{,}322} = \$75.962{,}96
\]

Interpretación: ¿Cuánto hubieras necesitado en 2018 para tener el mismo poder adquisitivo que esos \$100.000 pesos tienen en 2023? Respuesta: \$75.963. Lo que antes podías comprar con 76 mil pesos, ahora te cuesta \$100.000.

\subsection{Saldos de Dinero Real}

\begin{definicionbox}[Saldos de Dinero Real]
Los \textbf{saldos de dinero real} (o saldos reales) son la cantidad de dinero nominal dividida entre el nivel de precios. La demanda real del dinero se llama \textbf{Demanda por Saldos Reales}.
\end{definicionbox}

Cuanto más alto es el nivel de precios, más saldos nominales debe tener una persona para mantener el poder adquisitivo.

\subsection{Demanda por Saldos Reales}

La demanda por saldos reales depende de:
\begin{itemize}
    \item El \textbf{nivel de ingreso} ($Y$).
    \item La \textbf{tasa de interés} ($r$).
\end{itemize}

\subsubsection{Relación con la tasa de interés}

El tipo de interés es el \textbf{costo de oportunidad} de tener dinero. El costo de conservar el dinero es la tasa de interés que se pierde por tener dinero en lugar de otros activos.

\subsubsection{Relación con el nivel de ingreso}

Cuando la renta es alta, el gasto es elevado, por lo que la gente realiza más transacciones que exigen el uso del dinero. Esto viene de la \textbf{Teoría Cuantitativa del Dinero}:

\[
\underbrace{M}_{\text{Dinero}} \times \underbrace{V}_{\text{Velocidad}} = \underbrace{P}_{\text{Precio}} \times \underbrace{T}_{\text{Transacciones}}
\]

Ejemplo: $M = 100$ mil millones de pesos, $T = 400$ mil millones de transacciones, $P = 1{,}25$, $V = 5$ (cantidad de veces que debe cambiar de mano el dinero).

De aquí: $M \cdot V = P \cdot Y$, entonces $M = \dfrac{\barvar{P} \cdot Y}{\barvar{V}}$. Si $\uparrow Y \Rightarrow \uparrow M$.

\subsection{Forma Funcional de la Demanda por Saldos Reales}

\begin{resultadobox}[Función de Demanda por Saldos Reales]
\begin{equation}
    \left(\frac{M}{P}\right)^d = k \cdot Y - h \cdot r \tag{1}
\end{equation}
con $k, h > 0$.
\end{resultadobox}

Donde:
\begin{itemize}
    \item $k$ = Sensibilidad de la demanda por saldos reales al nivel de ingresos.
    \item $h$ = Sensibilidad de la demanda por saldos reales a la tasa de interés.
\end{itemize}

\subsubsection{Gráfico de la Demanda por Saldos Reales}

Puntos de corte: Si $r = 0$: $\left(\frac{M}{P}\right)^d = k \cdot Y$. Si $\left(\frac{M}{P}\right)^d = 0$: $r = \frac{k \cdot Y}{h}$.

\begin{center}
\begin{tikzpicture}
    \begin{axis}[
        axis lines=middle,
        xlabel={$\frac{M}{P}$}, ylabel={$r$},
        xmin=-0.5, xmax=8, ymin=-0.5, ymax=6,
        xtick=\empty, ytick=\empty,
        every axis x label/.style={at={(ticklabel* cs:1)}, anchor=west},
        every axis y label/.style={at={(ticklabel* cs:1)}, anchor=south},
        width=10cm, height=7cm,
    ]
    % Demanda por saldos reales para Y1
    \addplot[domain=0:5, thick, violet] {4 - 0.8*x};
    \node[below, violet] at (axis cs:5,0) {$k \cdot Y_1$};
    \node[left, violet] at (axis cs:0,4) {$\frac{k \cdot Y_1}{h}$};

    % Demanda para Y2 > Y1 (desplazada a la derecha)
    \addplot[domain=0:7, thick, red!70!black] {5.6 - 0.8*x};
    \node[below, red!70!black] at (axis cs:7,0) {$k \cdot Y_2$};

    % Flecha indicando aumento de Y
    \draw[->, thick, orange] (axis cs:3.5,1.5) -- (axis cs:5,1.5);
    \node[above, orange] at (axis cs:4.2,1.5) {$\uparrow Y$ ($Y_1 \to Y_2$)};

    % Tasa de interés fija
    \addplot[dashed, blue] coordinates {(0,2) (7,2)};
    \node[left, blue] at (axis cs:0,2) {$r$};

    % Intersecciones
    \addplot[only marks, mark=*, mark size=2pt, blue] coordinates {(2.5,2)};
    \addplot[only marks, mark=*, mark size=2pt, blue] coordinates {(4.5,2)};
    \addplot[dashed, blue!50] coordinates {(2.5,0) (2.5,2)};
    \addplot[dashed, blue!50] coordinates {(4.5,0) (4.5,2)};
    \node[below, blue] at (axis cs:2.5,0) {$Y_1$};
    \node[below, blue] at (axis cs:4.5,0) {$Y_2$};

    \end{axis}
\end{tikzpicture}
\end{center}

% ============================================================
\subsection{Oferta Monetaria}
% ============================================================

La oferta monetaria en Chile se clasifica en:
\begin{itemize}
    \item \textbf{M1:} Billetes, monedas en circulación, depósitos a la vista.
    \item \textbf{M2:} M1 + Depósitos a corto y mediano plazo.
    \item \textbf{M3:} M2 + Depósitos a largo plazo y otros activos líquidos.
\end{itemize}

¿Cómo influye el Banco Central de Chile (BCCh) en la oferta monetaria?
\begin{enumerate}
    \item Tasa de interés de referencia.
    \item Compra y venta de títulos de deuda.
    \item Regulación de la banca.
\end{enumerate}

\begin{resultadobox}[Oferta de Dinero]
\begin{equation}
    \left(\frac{M}{P}\right)^O = \frac{\barvar{M}}{\barvar{P}} \tag{2}
\end{equation}
La oferta de dinero real es exógena (fijada por el Banco Central).
\end{resultadobox}

\begin{center}
\begin{tikzpicture}
    \begin{axis}[
        axis lines=middle,
        xlabel={$M/P$}, ylabel={$r$},
        xmin=-0.5, xmax=8, ymin=-0.5, ymax=6,
        xtick=\empty, ytick=\empty,
        every axis x label/.style={at={(ticklabel* cs:1)}, anchor=west},
        every axis y label/.style={at={(ticklabel* cs:1)}, anchor=south},
        width=8cm, height=6cm,
    ]
    % Oferta monetaria (vertical)
    \addplot[thick, red!60!black] coordinates {(4,0) (4,5.5)};
    \node[above, red!60!black] at (axis cs:4,5.5) {$\left(\frac{M}{P}\right)^O$};
    \node[below] at (axis cs:4,0) {$\frac{\barvar{M}}{\barvar{P}}$};

    \end{axis}
\end{tikzpicture}
\end{center}

% ============================================================
\subsection{Equilibrio en el Mercado del Dinero}
% ============================================================

\begin{resultadobox}[Condición de Equilibrio del Mercado del Dinero]
\begin{equation}
    \left(\frac{M}{P}\right)^O = \left(\frac{M}{P}\right)^d \tag{3}
\end{equation}
Oferta = Demanda.
\end{resultadobox}

Reemplazando (1) y (2) en (3):
\begin{align}
    \frac{\barvar{M}}{\barvar{P}} &= k \cdot Y - h \cdot r \notag \\
    h \cdot r &= k \cdot Y - \frac{\barvar{M}}{\barvar{P}} \notag
\end{align}

\begin{resultadobox}[Curva LM]
\begin{equation}
    r = \frac{1}{h}\left(k \cdot Y - \frac{\barvar{M}}{\barvar{P}}\right) \tag{4}
\end{equation}
\end{resultadobox}

\begin{definicionbox}[Definición de la Curva LM]
La Curva LM muestra todas las combinaciones de tasa de interés y nivel de ingreso en donde el mercado del dinero está en equilibrio.
\end{definicionbox}

\subsubsection{Gráfico: Del Mercado del Dinero a la Curva LM}

\begin{center}
\begin{tikzpicture}
    % --- Mercado del Dinero (izquierda) ---
    \begin{axis}[
        name=dinero,
        axis lines=middle,
        xlabel={$M/P$}, ylabel={$r$},
        xmin=-0.5, xmax=8, ymin=-0.5, ymax=6,
        xtick=\empty, ytick=\empty,
        every axis x label/.style={at={(ticklabel* cs:1)}, anchor=west},
        every axis y label/.style={at={(ticklabel* cs:1)}, anchor=south},
        width=7cm, height=6cm,
        title={\small Mercado del Dinero},
    ]
    % Oferta
    \addplot[thick, red!60!black] coordinates {(4,0) (4,5.5)};

    % Demanda
    \addplot[domain=0:5, thick, violet] {4 - 0.8*x};

    % Equilibrio
    \addplot[only marks, mark=*, mark size=3pt, teal] coordinates {(4,0.8)};
    \addplot[dashed, gray] coordinates {(0,0.8) (4,0.8)};
    \node[left] at (axis cs:0,0.8) {$r^*$};
    \node[below] at (axis cs:4,0) {$\frac{\barvar{M}}{\barvar{P}}$};

    \end{axis}

    % --- Curva LM (derecha) ---
    \begin{axis}[
        at={(dinero.outer east)}, anchor=outer west, xshift=1.5cm,
        axis lines=middle,
        xlabel={$Y$}, ylabel={$r$},
        xmin=-0.5, xmax=8, ymin=-0.5, ymax=6,
        xtick=\empty, ytick=\empty,
        every axis x label/.style={at={(ticklabel* cs:1)}, anchor=west},
        every axis y label/.style={at={(ticklabel* cs:1)}, anchor=south},
        width=7cm, height=6cm,
        title={\small Curva LM},
    ]
    % Curva LM
    \addplot[domain=1:7, thick, violet] {-1.5 + 0.8*x};
    \node[above right, violet] at (axis cs:7,4.1) {$LM$};

    % Punto de equilibrio
    \addplot[only marks, mark=*, mark size=3pt, teal] coordinates {(5,2.5)};
    \addplot[dashed, gray] coordinates {(0,2.5) (5,2.5)};
    \addplot[dashed, gray] coordinates {(5,0) (5,2.5)};
    \node[left] at (axis cs:0,2.5) {$r^*$};
    \node[below] at (axis cs:5,0) {$Y_1$};
    \node[below] at (axis cs:1.875,0) {$\frac{\barvar{M}}{k\barvar{P}}$};

    \end{axis}
\end{tikzpicture}
\end{center}

Puntos de corte de la Curva LM:
\begin{itemize}
    \item Si $Y = 0$: $r = -\frac{\barvar{M}}{h\barvar{P}}$ (negativo).
    \item Si $r = 0$: $Y = \frac{\barvar{M}}{k\barvar{P}}$.
\end{itemize}

% ============================================================
\section{Equilibrio en el Mercado de Bienes y Mercado del Dinero}
% ============================================================

Tenemos dos ecuaciones y dos incógnitas ($Y$, $r$):

\begin{align}
    \text{Curva IS:} \quad &Y = \alpha_m(\barvar{A}_0 - b \cdot r) \tag{Curva IS} \\[6pt]
    \text{Curva LM:} \quad &r = \frac{1}{h}\left(k \cdot Y - \frac{\barvar{M}}{\barvar{P}}\right) \tag{Curva LM}
\end{align}

\subsection{Derivación de la Curva IS en función de $r$}

Primero ponemos la Curva IS en función de la tasa de interés ($r$):
\begin{align}
    Y &= \frac{\barvar{A}_0 - b \cdot r}{1 - c(1-t)} \notag \\
    (1 - c(1-t))Y &= \barvar{A}_0 - b \cdot r \notag \\
    b \cdot r &= \barvar{A}_0 - (1 - c(1-t))Y \notag
\end{align}

\begin{resultadobox}[Curva IS en función de $r$]
\begin{equation}
    r = \frac{\barvar{A}_0}{b} - \frac{Y}{\alpha_m \cdot b} \tag{5}
\end{equation}
\end{resultadobox}

Y la Curva LM:
\begin{resultadobox}[Curva LM]
\begin{equation}
    r = \frac{k \cdot Y}{h} - \frac{\barvar{M}}{h\barvar{P}} \tag{6}
\end{equation}
\end{resultadobox}

\subsection{Producción de Equilibrio ($Y^*$)}

Igualamos (5) y (6), $r = r$:
\[
\frac{\barvar{A}_0}{b} - \frac{Y}{\alpha_m b} = \frac{k \cdot Y}{h} - \frac{\barvar{M}}{h\barvar{P}}
\]

Despejamos $Y$:
\[
\frac{k \cdot Y}{h} + \frac{Y}{\alpha_m b} = \frac{\barvar{A}_0}{b} + \frac{\barvar{M}}{h\barvar{P}}
\]

\[
Y\left(\frac{k}{h} + \frac{1}{\alpha_m b}\right) = \frac{h\barvar{P}\barvar{A}_0 + b\barvar{M}}{bh\barvar{P}}
\]

\[
Y\left(\frac{\alpha_m b k + h}{h \alpha_m b}\right) = \frac{h\barvar{P}\barvar{A}_0 + b\barvar{M}}{bh\barvar{P}}
\]

Simplificando y definiendo:
\[
\gamma = \frac{\alpha_m}{1 + \dfrac{kb\alpha_m}{h}}
\]

\begin{resultadobox}[Producción de Equilibrio IS-LM]
\begin{equation}
    Y^* = \gamma \cdot \barvar{A}_0 + \gamma \cdot \frac{b}{h} \cdot \frac{\barvar{M}}{\barvar{P}} \tag{7}
\end{equation}
\end{resultadobox}

Este resultado es equivalente a la ecuación (8a) de la página 243 del Dornbusch:
\[
Y = \gamma \barvar{A} + \gamma \frac{b}{h}\frac{\barvar{M}}{\barvar{P}}
\]

Donde:
\begin{itemize}
    \item $\barvar{A} = \barvar{A}_0$ (Gasto Autónomo).
    \item $\gamma = \dfrac{\alpha_G}{1 + \frac{k \alpha_G b}{h}}$, con $\alpha_G = \alpha_m$ (multiplicador del gasto).
\end{itemize}

\subsection{Tasa de Interés de Equilibrio ($r^*$)}

Para encontrar $r^*$, reemplazamos (7) en la Curva LM (6):
\begin{align}
    r &= \frac{k \cdot Y^*}{h} - \frac{\barvar{M}}{h\barvar{P}} \notag \\[6pt]
    r &= \frac{k}{h}\left(\gamma \barvar{A}_0 + \gamma \frac{b}{h}\frac{\barvar{M}}{\barvar{P}}\right) - \frac{\barvar{M}}{h\barvar{P}} \notag
\end{align}

Desarrollando y simplificando:

\begin{resultadobox}[Tasa de Interés de Equilibrio IS-LM]
\begin{equation}
    r^* = \gamma \cdot \frac{k}{h} \cdot \barvar{A}_0 - \gamma \cdot \frac{1}{h \alpha_m} \cdot \frac{\barvar{M}}{\barvar{P}} \tag{8}
\end{equation}
\end{resultadobox}

Este resultado es equivalente a la ecuación (9a) de la página 244 del Dornbusch:
\[
i = \gamma \frac{k}{h}\barvar{A} - \gamma \frac{1}{h\alpha_G}\frac{\barvar{M}}{\barvar{P}}
\]
donde $i = r$ (en este modelo sin inflación, la tasa nominal es igual a la real).

\end{document}
